\documentclass[12pt,a4paper]{article}

% ABNT formatting packages
\usepackage[inner=30mm,outer=20mm,top=30mm,bottom=20mm]{geometry}
\usepackage[utf8]{inputenc}
\usepackage[T1]{fontenc}
\usepackage[brazilian]{babel}
\usepackage{graphicx}
\usepackage{float}
\usepackage[colorlinks=true,linkcolor=black,urlcolor=black,citecolor=black]{hyperref}
\usepackage{indentfirst}
\usepackage{setspace}
\usepackage{helvet}
\usepackage[font=normalsize,labelfont=bf,labelsep=endash,justification=centering,position=above]{caption}

\renewcommand{\familydefault}{\sfdefault}
\onehalfspacing
\setlength{\parindent}{1.5cm}

\input{../../../out/drive-links.tex}

\begin{document}

\section{Diagrama de Implantação}

Na Figura~\ref{fig:implantacao} é apresentado o diagrama de implantação do sistema,
descrevendo a arquitetura física e a distribuição dos componentes nos nós de infraestrutura.

\begin{figure}[H]
    \caption{Diagrama de implantação do sistema de reserva de salas}
    \label{fig:implantacao}
    \centering
    \includegraphics[width=0.95\textwidth,height=0.75\textheight,keepaspectratio]{../../../out/src/plantuml/deployment.png}

    {\footnotesize \textbf{Fonte:} O autor.
    \textbf{Acesso:} \href{\linkDeployment}{Diagrama de Implantação}.}
\end{figure}

\end{document}
